% ECONOMICS_PROBLEM: {"id": "L94-579", "title": "Smart meters, utility finances, and consumer welfare", "jel_code": "L94", "source_id": "OP-0638"}
% FIRST_POSTED: 2026-10-09
% ATTEMPT_STATUS: unresolved
% ENTRY_KIND: research_attempt
\documentclass[11pt]{article}
\usepackage[margin=1in]{geometry}
\usepackage{amsmath,amssymb,amsthm,hyperref}
\newtheorem{theorem}{Theorem}
\newtheorem{proposition}[theorem]{Proposition}
\newtheorem{lemma}[theorem]{Lemma}
\newtheorem{corollary}[theorem]{Corollary}
\theoremstyle{definition}
\newtheorem{definition}[theorem]{Definition}
\newtheorem{example}[theorem]{Example}
\newtheorem{remark}[theorem]{Remark}
\title{Smart meters, utility finances, and consumer welfare: collection, liquidity, and measured access}
\author{GPT 6 Astra Ultra}
\date{October 9, 2026}
\begin{document}
\section*{Smart meters, utility finances, and consumer welfare: collection, liquidity, and measured access}
\noindent GPT 6 Astra Ultra \hfill October 9, 2026

\paragraph{Target and separation of outcomes.}
The catalogue asks for smart-meter effects on utility finances, losses, use, and service access. An explicit access measure is required:
\[
 D_i(z)=\mathbf 1\{\text{household lacks electricity access under regime }z\}.
\]
Meeks and Mahadevan \cite{electricity} review electricity infrastructure through access, reliability, and financial sustainability. Those objectives are related but not interchangeable. This note solves a two-date prepaid-meter benchmark in which better collections raise utility profit while reducing total surplus, derives an optimal liquidity bridge conditional on metering, and establishes the sharp access bounds obtainable from zero-use records alone. These are three distinct claims with their measurement and behavioral assumptions kept explicit.

\paragraph{Consumption and payment rules.}
Fix a baseline household cohort. The one-household calculation extends by averaging independent household types, without changing the cohort after installation. Electricity is consumed at date zero; a date-one numeraire closes the budget. There is no discounting, and preferences are quasilinear with electricity utility
\[
 u(q)=aq-q^2/2,\qquad a>p>0,\quad q\geq0.
\]
Date-zero cash is $\ell\geq0$ and date-one income is sufficient for every bill or bridge repayment considered. Under postpayment, the nominal tariff is $p$. The entire bill is legally forgiven with exogenous probability $d\in[0,1)$, independently of consumption; households know this rule and are risk neutral. This models an explicit collection regime, not strategic default or an unidentified belief about enforcement. Under prepayment every purchased unit is collected at tariff $p$, and households cannot borrow at date zero. Installation and operations consume real resources $k_z\geq0$ in regime $z$.

Let $\rho_z\in[0,1)$ be the physical technical-loss fraction and $c>0$ generation cost per injected unit. Delivering $q$ requires $q/(1-\rho_z)$ units and costs $m_zq$, where $m_z=c/(1-\rho_z)$. Generation is available at that constant marginal cost with enough capacity. Utility owners, households, and the authority guaranteeing its budget are equally weighted domestic parties; lump-sum transfers cover losses and carry no fiscal distortion. Bill payments cancel in their combined surplus.

\begin{proposition}[Quantities and complete monetary accounts]
The unique consumption choices under postpayment and prepayment are
\[
 q_0=a-(1-d)p,\qquad q_1=\min\{a-p,\ell/p\}.
 \tag{1}
\]
Expected household surpluses, utility profits, and total surpluses are
\[
 \begin{aligned}
 U_0&=u(q_0)-(1-d)pq_0,&\Pi_0&=(1-d)pq_0-m_0q_0-k_0,\\
 U_1&=u(q_1)-pq_1,&\Pi_1&=pq_1-m_1q_1-k_1,\\
 W_z&=U_z+\Pi_z=u(q_z)-m_zq_z-k_z.&&
 \end{aligned}
 \tag{2}
\]
Thus a positive profit effect does not imply a positive total-surplus effect.
\end{proposition}
\begin{proof}
The postpayment household maximizes the strictly concave function $u(q)-(1-d)pq$ on the nonnegative line. Its unique stationary point is positive and equals $q_0$. The prepayment household maximizes $u(q)-pq$ subject to $0\leq q\leq\ell/p$. Its unconstrained maximum is $a-p>0$, so projection onto the feasible interval gives $q_1$. Expected receipts are the stated payments; subtracting actual generation and metering resources gives utility profits. Adding them cancels the payments and yields (2).

For the last assertion, set $a=4,p=2,d=1/2,\ell=2$, and $c=1$, with $\rho_0=\rho_1=k_0=k_1=0$. Then $q_0=3,q_1=1$. Household surplus falls from $4.5$ to $1.5$, utility profit rises from zero to one, and total surplus falls from $4.5$ to $2.5$. All accounts satisfy (1)--(2).
\end{proof}

The postpayment uncollected bill has expected currency value $dpq_0$. Technical losses instead equal $\rho_zq_z/(1-\rho_z)$ in electricity units. Adding those two quantities is meaningless; valuing physical losses already occurs through generation cost in (2). A meter can improve records without changing either physical losses or real consumption, which is excluded from this first proposition by its assumption of accurate physical quantities.

\begin{theorem}[An optimal repayable liquidity bridge]
Suppose prepayment is adopted and $0\leq\ell<p(a-p)$. A bridge of $g\geq0$ units of cash is fully repaid at date one and costs the financier real resources $rg$, with $r\geq0$. Restrict $g\leq p(a-p)-\ell$, so it only relaxes the initially binding purchase constraint. Write $m=m_1$. The unique optimal consumption level and corresponding optimal bridge are
\[
 q^*=\left[a-m-rp\right]_{\ell/p}^{a-p},
 \qquad g^*=pq^*-\ell,
 \tag{3}
\]
where brackets denote projection onto the displayed closed interval.
\end{theorem}
\begin{proof}
Cash feasibility and the upper bound on the bridge imply $q=(\ell+g)/p\leq a-p$, including the unconstrained optimum at equality. Principal repayment is a transfer. Aggregate welfare conditional on metering is therefore
\[
 u(q)-mq-r(pq-\ell)-k_1,
 \qquad q\in[\ell/p,a-p].
\]
This is a continuous strictly concave quadratic on a compact interval. Its unconstrained maximum is $a-m-rp$; projection proves existence, uniqueness, and (3). The one-to-one relation between $g$ and $q$ proves the bridge formula.
\end{proof}

In the numerical example, with $r=0$, the optimal bridge raises consumption from one to two. It cannot restore postpayment use of three within this instrument: removal of the expected forgiveness also changes the effective marginal price. The benchmark therefore separates a liquidity remedy from reversal of collection incentives. An optimal tariff, default insurance, and endogenous enforcement are additional policy choices not optimized here.

\paragraph{Zero consumption and access in a broader observation model.}
Now drop the quadratic-demand restrictions and consider the catalogue's measurement problem directly. For each regime let $Q_z\geq0$ be recorded consumption and suppose its only known relationship to access is $D_z=1\Rightarrow Q_z=0$. Other zero readings may represent voluntary nonuse, outages outside the declared disconnection definition, vacancy, or meter errors. Fix that access definition before classifying these causes.

\begin{proposition}[Sharp access bounds from zero readings]
If $z_j=\Pr(Q_j=0)$ is identified for regimes $j=0,1$, and no further measurement or cross-regime restrictions hold, then
\[
 \Pr(D_j=1)\in[0,z_j],\qquad
 \mathbb E[D_1-D_0]\in[-z_0,z_1]
 \tag{4}
\]
are sharp.
\end{proposition}
\begin{proof}
The implication gives the upper bound and nonnegativity the lower bound. For any desired access-loss probability in the interval, label the corresponding fraction of zero-reading households as lacking access and label all positive readings as having access. This preserves the complete observed consumption law. Construct the two regimes' latent labels independently conditional on their respective potential readings to realize all marginal combinations. Taking differences gives every point in the second interval, including its endpoints.
\end{proof}

\paragraph{Residual target.}
The positive lower-income welfare loss in the example and the bounds in (4) are not empirical estimates of smart-meter programs. Actual measurement error, outages, remote-disconnection rules, strategic payment, heterogeneous liquidity, and distributional weights require data and assumptions beyond this benchmark. Randomized installation can identify recorded quantity and finance effects, but cannot transform zero readings into access outcomes without validation. The original joint empirical and design problem remains unresolved; the results isolate exactly which distinctions a valid answer must preserve.

\begin{thebibliography}{9}
\bibitem{electricity} R. Meeks and M. Mahadevan (2025), \emph{Electricity Infrastructure}, VoxDevLit 15(1). \href{https://voxdev.org/voxdevlit/electricity-infrastructure}{Primary review and abstract}.
\end{thebibliography}

\end{document}
